% !TEX TS-program = pdflatex
% !TEX encoding = UTF-8 Unicode

% This is a simple template for a LaTeX document using the "article" class.
% See "book", "report", "letter" for other types of document.

\documentclass[12pt]{report} % use larger type; default would be 10pt

\usepackage[utf8]{inputenc} % set input encoding (not needed with XeLaTeX)

%%% Examples of Article customizations
% These packages are optional, depending whether you want the features they provide.
% See the LaTeX Companion or other references for full information.

%%% PAGE DIMENSIONS
\usepackage{geometry} % to change the page dimensions
\geometry{a4paper} % or letterpaper (US) or a5paper or....
% \geometry{margin=2in} % for example, change the margins to 2 inches all round
% \geometry{landscape} % set up the page for landscape
%   read geometry.pdf for detailed page layout information

\usepackage{graphicx} % support the \includegraphics command and options

% \usepackage[parfill]{parskip} % Activate to begin paragraphs with an empty line rather than an indent

%%% PACKAGES
\usepackage{booktabs} % for much better looking tables
\usepackage{array} % for better arrays (eg matrices) in maths
\usepackage{paralist} % very flexible & customisable lists (eg. enumerate/itemize, etc.)
\usepackage{verbatim} % adds environment for commenting out blocks of text & for better verbatim
\usepackage{subfig} % make it possible to include more than one captioned figure/table in a single float
\usepackage[final]{pdfpages}
% These packages are all incorporated in the memoir class to one degree or another...

%%% HEADERS & FOOTERS
\usepackage{fancyhdr} % This should be set AFTER setting up the page geometry
\pagestyle{fancy} % options: empty , plain , fancy
\renewcommand{\headrulewidth}{0pt} % customise the layout...
\lhead{}\chead{}\rhead{}
\lfoot{}\cfoot{\thepage}\rfoot{}

%%% SECTION TITLE APPEARANCE
\usepackage{sectsty}
\allsectionsfont{\sffamily\mdseries\upshape} % (See the fntguide.pdf for font help)
% (This matches ConTeXt defaults)

%%% RULE

\newcommand{\HRule}{\rule{\linewidth}{0.5mm}}

%%% BIBLIOGRAPHY

\usepackage{apacite}                           %bibliography in apa-style

%%% ToC (table of contents) APPEARANCE
\usepackage[nottoc,notlof,notlot]{tocbibind} % Put the bibliography in the ToC
\usepackage[titles,subfigure]{tocloft} % Alter the style of the Table of Contents
\renewcommand{\cftsecfont}{\rmfamily\mdseries\upshape}
\renewcommand{\cftsecpagefont}{\rmfamily\mdseries\upshape} % No bold!

\setcounter{secnumdepth}{-2}

%%% TABLES

\renewcommand{\arraystretch}{1.2}

%%% END Article customizations

%%% The "real" document content comes below...

\begin{document}

\begin{titlepage}

\begin{center}


% Upper part of the page
\includegraphics[width=1\textwidth]{./logo}\\[1cm]    

\textsc{\Large Bachelor Thesis}\\[0.5cm]
\textsc{\Large {[}201000166{]}}\\[0.5cm]


% Title
\HRule \\[0.4cm]
{ \huge \bfseries Teaching quantum mechanics using qCraft}\\[0.4cm]

\HRule \\[1.5cm]

% Author and supervisor
\begin{minipage}{0.4\textwidth}
\begin{flushleft} \large
\emph{Author:}\\
Micha \textsc{van den Enk} \\
{[}s1004654{]} \\
\end{flushleft}
\end{minipage}
\begin{minipage}{0.4\textwidth}
\begin{flushright} \large
\emph{Supervisors:} \\
Dr. H. H. \textsc{Leemkuil} \\
Second \textsc{supervisor} \\
\end{flushright}
\end{minipage}

\vfill

% Bottom of the page
{\large \today}

\end{center}

\end{titlepage}

\setcounter{tocdepth}{1}
\tableofcontents

\chapter{Preface}

%%% ANALYSIS

\part{Analyses}



% Context Analysis

\chapter{Context Analysis}

%needs assessment

\section{Needs Assessment}

\subsection{The condition}

\subsection{The problem}

\subsection{Cause of the problem}

\subsection{Is the solution learning}

\subsection{Instruction currently offered}

\subsection{Nature of the innovation}

\subsection{Learning goals}

\subsection{Priority and suitability}

%learning environment

\section{Learning Environment}

\subsection{Teachers}

\subsection{Existing curricula}

\subsection{Equipment}

\subsection{Facilities}

\subsection{Organization}

\subsection{Larger system}

% Learner analysis

\chapter{Learner Analysis}

%stable similarities

\section{Stable Similarities}

\subsection{Sensory capacities}

\subsection{Information processing}

\subsection{Types and conditions of learning}

%changing similarities

\section{Changing Similarities}

\subsection{Intellectual development processes}

\subsection{Language development processes}

\subsection{Psychosocial development processes}

\subsection{Moral development processes}

\subsection{Other development processes}

%stable differences

\section{Stable Differences}

\subsection{Aptitudes}

\subsection{Cognitive styles}

\subsection{Psychosocial traits}

\subsection{Gender, ethnicity, \& racial group}

%changing differences

\section{Changing Differences}

\subsection{Intellectual development state}

\subsection{Other development state}

\subsection{General prior learning}

\subsection{Specific prior learning}

% Task analysis

\chapter{Task Analysis}

\section{Learning goal}

\section{Types of learning}

\section{Information-processing analysis}

\section{Prerequisite analysis}

\section{Learning objectives}

\section{Test specifications}

% Theoretic Framework

\chapter{Theoretic Framework}

%%% DESIGN

\part{Design}

%%% DEVELOPMENT

\part{Development}

%%% FORMATIVE EVALUATION

\part{Formative Evaluation}

\bibliographystyle{apacite}
\bibliography{references}

\end{document}
